\section{Disscusion}
% \myparagraph{Limitation.}
% first latent diffusion model in motion,
As the first trial, to our best knowledge, exploring human motion generation using language models, the proposed MotionGPT still owns limitations as follows.
%
MotionGPT only utilizes motion on articulated human bodies, while many other works focus on faces~\cite{karras2017audio, cao2018sparse}, hands~\cite{romero2022embodied, li2022nimble, li2021piano} and even animal~\cite{Rueeg:CVPR:2022, Zuffi:CVPR:2018} motion.
%
Besides, our method is also restricted to multiple humans without modeling human-object, or human-environment interactions~\cite{shafir2023priormdm}.
%
It is interesting to model the human interaction scenarios in a motion-language framework and generate controllable motions~\cite{shafir2023priormdm}.


We summarize the proposed MotionGPT as a uniform motion-language framework to generate plausible human motion and natural language descriptions through prompt-based instructions. Compared to the compatible motion diffusion methods~\cite{chen2023mld, mdm2022human}, our MotionGPT produces competitive results on motion generation, motion captioning, motion prediction, and motion in-between using only one pre-trained generative model. With the advancement of large language data and models~\cite{raffel2020t5,chung2022flant5}, MotionGPT is also capable of addressing natural question-to-answer tasks.
%
Extensive experiments on various human motion-relevant tasks demonstrate the effectiveness and extendibility of MotionGPT.

% We present a robust template-based approach to capture challenging 3D human motions using only a single RGB camera, which is in a novel two-stage learning-and-optimization framework to make full usage of multi-modal references.
% %
% Our hybrid motion inference learns the challenging motion details from various supervised references modalities, while our robust motion optimization further improves the tracking accuracy by extracting the reliable motion hints in the input image reference.
% %
% Our experimental results demonstrate the effectiveness and robustness of ChallenCap in capturing challenging human motions in various scenarios.
% %
% We believe that it is a significant step to enable robust 3D capture of challenging human motions,
% with many potential applications in VR/AR, and performance evaluation for gymnastics, sports, and dancing.

% which compares favorably to the state-of-the-arts.
%it suffers from complex motion patterns and severe self-occlusion under the monocular setting. 
